\documentclass[../main.tex]{subfiles}

\begin{document}

\chapter*{Landasan Teori dan Peta Konsep}
\addcontentsline{toc}{chapter}{Landasan Teori dan Peta Konsep}

Materi buku ini bertumpu pada tiga hal: konsep dan definisi kunci, teori utama
kriptografi, serta peta keterkaitan antar-bab. Bab 1 sampai Bab 3 berperan
sebagai landasan teori dan konsep dasar; bab-bab berikutnya adalah unit materi
inti yang menerapkan landasan itu.

\section*{Konsep dan Definisi Kunci}

Definisi istilah berikut dipakai berulang di seluruh buku (definisi formalnya
diberikan pada bab yang ditunjuk).

\begin{description}
    \item[Kriptografi] Ilmu dan seni menjaga keamanan pesan, mencakup
          enkripsi, dekripsi, serta layanan keamanan lain (Bab 1).
    \item[Plainteks dan cipherteks] Pesan asli dan pesan hasil penyandian
          (Bab 1).
    \item[Kunci] Informasi rahasia yang mengendalikan proses enkripsi dan
          dekripsi; keamanan sistem bergantung pada kerahasiaan kunci, bukan
          pada kerahasiaan algoritma (Bab 1 dan Bab 2).
    \item[Layanan keamanan] Kerahasiaan, integritas, ketersediaan,
          otentikasi, dan nirsangkal (Bab 2).
    \item[Aritmetika modulo dan invers modulo] Operasi sisa pembagian beserta
          balikannya, landasan matematis hampir semua algoritma pada buku ini
          (Bab 3).
    \item[Kriptanalisis] Upaya memperoleh plainteks atau kunci tanpa hak
          (Bab 2).
\end{description}

\section*{Teori Utama}

\begin{itemize}
    \item \textbf{Keamanan melalui kerahasiaan kunci.} Sesuai prinsip Kerckhoffs,
          algoritma boleh diketahui umum; yang harus dirahasiakan hanya kuncinya.
    \item \textbf{Kerahasiaan sempurna.} \textit{One-time pad} mencapai keamanan
          teoretis bila kunci acak, sepanjang pesan, dan tidak pernah dipakai
          ulang (Bab 4).
    \item \textbf{Konfusi dan difusi.} Shannon merumuskan dua sifat yang harus
          dimiliki block cipher yang kuat (Bab 6).
    \item \textbf{Jaringan Feistel.} Kerangka yang memungkinkan dekripsi memakai
          mesin yang sama dengan enkripsi, dan mendasari DES (Bab 6 dan Bab 7).
    \item \textbf{Kesulitan komputasi.} Keamanan kriptografi kunci-publik
          bersandar pada persoalan yang sulit dipecahkan---faktorisasi bilangan
          besar (RSA), logaritma diskret (Diffie--Hellman, ElGamal), dan
          logaritma diskret pada kurva eliptik (ECC) (Bab 10--14).
    \item \textbf{Fungsi satu arah dan tanda tangan digital.} Integritas dan
          nirsangkal dibangun dari fungsi hash dan kunci privat penandatangan
          (Bab 15 dan Bab 16).
\end{itemize}

\section*{Peta Konsep dan Struktur Materi}

Bagan berikut merangkum alur materi buku: dari landasan, melalui kriptografi
klasik dan modern, hingga mekanisme integritas dan otentikasi.

\begin{center}
\begin{tikzpicture}[
    font=\scriptsize,
    kotak/.style={draw, rounded corners=2pt, align=center, inner sep=4pt,
                  text width=2.35cm, minimum height=1.15cm, fill=blue!5},
    akar/.style={draw, rounded corners=3pt, align=center, inner sep=5pt,
                  font=\small\bfseries, fill=blue!20},
    panah/.style={-{Stealth[length=2mm]}, thick, blue!55!black}
]
\node[akar] (akar) {Kriptografi};
\node[kotak, below=0.9cm of akar, xshift=-6.0cm] (b1)
    {\textbf{Fondasi}\\{\tiny Bab 1--3}};
\node[kotak, below=0.9cm of akar, xshift=-3.0cm] (b2)
    {\textbf{Klasik}\\{\tiny Bab 4}};
\node[kotak, below=0.9cm of akar] (b3)
    {\textbf{Simetris Modern}\\{\tiny Bab 5--9}};
\node[kotak, below=0.9cm of akar, xshift=3.0cm] (b4)
    {\textbf{Asimetris}\\{\tiny Bab 10--14}};
\node[kotak, below=0.9cm of akar, xshift=6.0cm] (b5)
    {\textbf{Integritas \& Otentikasi}\\{\tiny Bab 15--16}};
\foreach \b in {b1,b2,b3,b4,b5} \draw[panah] (akar) -- (\b);
\end{tikzpicture}
\end{center}

Kelima cabang pada bagan di atas adalah lima bagian materi buku.

\begin{enumerate}[label=\textbf{Bagian \arabic*.}, leftmargin=*]
    \item \textbf{Fondasi (Bab 1--3).} Pengantar dan urgensi kriptografi,
          layanan keamanan dan serangan, serta landasan matematika. Menopang
          Sub-CPMK 2.1--2.3.
    \item \textbf{Kriptografi klasik (Bab 4).} Cipher substitusi dan
          transposisi, abjad-majemuk (Vigen\`ere), polygram (Playfair, Hill),
          Enigma, Affine, dan \textit{one-time pad}. Menopang Sub-CPMK 1.1.
    \item \textbf{Kriptografi simetris modern (Bab 5--9).} Stream cipher, block
          cipher dan mode operasinya, DES, AES, serta block cipher lain (GOST,
          RC5/RC6, Blowfish, Twofish). Menopang Sub-CPMK 1.2.
    \item \textbf{Kriptografi asimetris (Bab 10--14).} RSA, pertukaran kunci
          Diffie--Hellman, ElGamal, kripto knapsack, dan kurva eliptik (ECC).
          Menopang Sub-CPMK 1.3.
    \item \textbf{Integritas dan otentikasi (Bab 15--16).} Fungsi hash dan MAC,
          tanda tangan digital, PKI, serta protokol SSL/TLS. Menopang
          Sub-CPMK 1.4.
\end{enumerate}

\end{document}
